\documentclass[12pt]{ximera}
\title{Exam 2 (2022)}
\begin{document}
\begin{abstract}
An archived exam on Chapter 2 and the start of Chapter 3: gridlock, anarchy, and veto power, Banzhaf power, a mayor's Shapley-Shubik pivots, counting sequential coalitions, continuous versus discrete assets, and divider-chooser.
\end{abstract}
\maketitle

\section*{Exam 2 (2022)}

This exam allowed one handwritten sheet of notes.

\begin{problem}
Consider the following weighted voting systems.
\[
\begin{array}{ll}
1) & $[26:10,6,5,4]$\\
2) & $[35:9,8,7,6]$\\
3) & $[13:11,9,6,2]$\\
4) & $[14:10,9,6,4]$\\
5) & $[11:7,4,2,1]$\\
\end{array}
\]

\begin{prompt}
\textbf{(a)} Which systems have gridlock?
\begin{selectAll}
\choice[correct]{System 1}
\choice[correct]{System 2}
\choice{System 3}
\choice{System 4}
\choice{System 5}
\end{selectAll}
\begin{feedback}
Gridlock means the quota exceeds the total weight, so no coalition can win.
System 1: $10+6+5+4=25<26$. System 2: $9+8+7+6=30<35$.
\end{feedback}
\end{prompt}

\begin{prompt}
\textbf{(b)} Which systems have anarchy?
\begin{selectAll}
\choice{System 1}
\choice{System 2}
\choice[correct]{System 3}
\choice[correct]{System 4}
\choice{System 5}
\end{selectAll}
\begin{feedback}
Anarchy means the quota is at most half the total weight, so two disjoint coalitions
could both win. System 3: $11+2\ge13$ and $9+6\ge13$. System 4: $10+4\ge14$ and $9+6\ge14$. System 5 is neither.
\end{feedback}
\end{prompt}

\begin{prompt}
\textbf{(c)} Which players in system 5 have veto power?
\begin{selectAll}
\choice[correct]{$P_1$}
\choice[correct]{$P_2$}
\choice{$P_3$}
\choice{$P_4$}
\end{selectAll}
\begin{feedback}
Without $P_1$ the others total $4+2+1=7<11$; without $P_2$ they total $7+2+1=10<11$. Without $P_3$ or $P_4$ the rest still reach $11$.
\end{feedback}
\end{prompt}
\end{problem}

\begin{problem}
Four business partners make decisions with the weighted voting system $[12:7,5,3,2]$, with
players $P_1,\dots,P_4$ in that order.

\begin{prompt}
\textbf{(a)} How many winning coalitions are there? $\answer{5}$
\begin{feedback}
$\set{P_1,P_2}=12$, $\set{P_1,P_2,P_3}=15$, $\set{P_1,P_2,P_4}=14$, $\set{P_1,P_3,P_4}=12$, and the grand coalition $17$. ($\set{P_2,P_3,P_4}=10$ loses.)
\end{feedback}
\end{prompt}

\begin{prompt}
\textbf{(b)} Find the critical counts and the Banzhaf power distribution, as fractions.

Critical counts: $P_1$: $\answer{5}$, $P_2$: $\answer{3}$,
$P_3$: $\answer{1}$, $P_4$: $\answer{1}$

\smallskip
$\beta_1=\answer{1/2}$, $\beta_2=\answer{3/10}$,
$\beta_3=\answer{1/10}$, $\beta_4=\answer{1/10}$
\begin{feedback}
$P_1$ is critical in all five. $P_2$ is critical in $\set{P_1,P_2}$, $\set{P_1,P_2,P_3}$, and $\set{P_1,P_2,P_4}$. $P_3$ and $P_4$ are critical only in $\set{P_1,P_3,P_4}$. The total is $10$, so $\frac5{10},\frac3{10},\frac1{10},\frac1{10}$.
\end{feedback}
\end{prompt}
\end{problem}

\begin{problem}
A city council has $9$ seats: $8$ councilors and the mayor, $M$. They use simple majority ($5$ votes), but a No vote by the mayor can only be overridden by $6$ or more Yes votes. Write
$\underline{\ }$ for a councilor.

\begin{prompt}
\textbf{(a)} In a sequential coalition with $M$ in slot $6$, which slot holds the
pivotal player? $\answer{6}$
\begin{feedback}
The first $5$ players are councilors: a majority, but not enough to override the mayor. Adding $M$ makes a majority that includes the mayor, so $M$ in slot $6$ is pivotal.
\end{feedback}
\end{prompt}

\begin{prompt}
\textbf{(b)} In a sequential coalition with $M$ in slot $4$, which slot
holds the pivotal player? $\answer{5}$
\begin{feedback}
The first $5$ players include the mayor and form a majority, so they win; the councilor in slot $5$ is pivotal.
\end{feedback}
\end{prompt}

\begin{prompt}
\textbf{(c)} In which slots is the mayor pivotal?
\begin{selectAll}
\choice{1}
\choice{2}
\choice{3}
\choice{4}
\choice[correct]{5}
\choice[correct]{6}
\choice{7}
\choice{8}
\choice{9}
\end{selectAll}
\begin{feedback}
In slots $1$--$5$, the coalition first wins at slot $5$ (a majority with the mayor), so $M$ is pivotal only in slot $5$. In slots $7$--$9$, six councilors win without $M$ at slot $6$. So $M$ is pivotal in slots $5$ and $6$: a Shapley-Shubik power of $\frac29$.
\end{feedback}
\end{prompt}
\end{problem}

\begin{problem}
How many ways are there to fill slots $1,2,3,5,7,8,9$ (seven slots) of a sequential coalition
with seven different councilor names A, B, C, D, E, F, G? (This is unrelated to the
previous problem.) $\answer{5040}$
\begin{feedback}
Only the number of slots and names matters: $7!=7\cdot6\cdot5\cdot4\cdot3\cdot2\cdot1=5040$.
\end{feedback}
\end{problem}

\begin{problem}
Each of the following is to be divided among $5$ players. Is it a continuous or a
discrete fair-division problem?

\begin{prompt}
\textbf{(a)} A pizza with one part mozzarella-prosciutto, one part meat lover's, one part five-cheese, and one part mushroom-artichoke.
\begin{multipleChoice}
\choice[correct]{Continuous}
\choice{Discrete}
\end{multipleChoice}
\begin{feedback}
Continuous: a pizza can be cut anywhere.
\end{feedback}
\end{prompt}
\begin{prompt}
\textbf{(b)} A silver locket that is impossible to open, a cracked seal made of unobtanium, and a golden cup with a house crest.
\begin{multipleChoice}
\choice{Continuous}
\choice[correct]{Discrete}
\end{multipleChoice}
\begin{feedback}
Discrete: these are indivisible objects.
\end{feedback}
\end{prompt}
\begin{prompt}
\textbf{(c)} $200{,}000$ barrels of light crude, $75{,}000{,}000$ gallons of liquefied natural gas, and $50{,}300$ metric tonnes of anthracite coal.
\begin{multipleChoice}
\choice[correct]{Continuous}
\choice{Discrete}
\end{multipleChoice}
\begin{feedback}
Continuous: bulk quantities can be split in any proportion.
\end{feedback}
\end{prompt}
\end{problem}

\begin{problem}
Charli and Denise have inherited a house and a plot of land. They both think the house is
worth \$160K. Charli thinks the land is worth \$120K, and Denise thinks it is
worth \$200K. The land can be divided in any way necessary. Denise is the divider and
makes share $s_1$ from the house and some of the land, and share $s_2$ from the rest of the
land.

\begin{prompt}
\textbf{(a)} What fraction of the land goes with the house in $s_1$, so that both shares
are equal in Denise's opinion? $\answer{1/10}$
\begin{feedback}
Denise values everything at $160+200=360$ (in thousands), so each share must be worth
$180$ to her. The house needs another $20$ in land, which is $\frac{20}{200}=\frac1{10}$
of the land. So $s_1=H+\frac1{10}L$ and $s_2=\frac9{10}L$.
\end{feedback}
\end{prompt}

\begin{prompt}
\textbf{(b)} How much is each share worth in Charli's opinion, in thousands of
dollars?

$s_1$: $\answer{172}$ \quad $s_2$: $\answer{108}$
\begin{feedback}
$s_1=160+\frac1{10}\cdot120=172$ and $s_2=\frac9{10}\cdot120=108$.
\end{feedback}
\end{prompt}

\begin{prompt}
\textbf{(c)} Which share will Charli choose?
\begin{multipleChoice}
\choice[correct]{$s_1$}
\choice{$s_2$}
\end{multipleChoice}
\begin{feedback}
$s_1$, worth $172>140$, half of Charli's total of $280$. (Denise, the divider, gets $s_2$, worth exactly $180$ to her.)
\end{feedback}
\end{prompt}
\end{problem}

\end{document}
